\section*{APPENDIX}
\phantomsection
\addcontentsline{toc}{section}{APPENDIX}
\setcounter{table}{0}
\setcounter{figure}{0}
\setcounter{equation}{0}
\renewcommand{\thetable}{A.\arabic{table}}
\renewcommand{\thefigure}{A.\arabic{figure}}
\renewcommand{\theequation}{A.\arabic{equation}}
\renewcommand{\theHtable}{A.\arabic{table}}
\renewcommand{\theHfigure}{A.\arabic{figure}}
\renewcommand{\theHequation}{A.\arabic{equation}}

This appendix documents the MATLAB implementation used to generate the numerical results in Chapter 5. The purpose is not to repeat the theoretical derivations or the simulation tables already presented in the main chapters, but to provide a compact reference for the source-code organization and the reproduction workflow.

\subsection*{Appendix A. MATLAB Source-Code Organization}

The simulation framework is organized according to the same technical structure used in the thesis. The baseline OTFS files implement the transmitter, channel, demodulator, and detector described in Chapter 3. The OTFS-IM files implement the index-modulation mapping and pattern-selection procedures developed in Chapter 4. The simulation and reporting scripts generate the BER, error-decomposition, pattern-validation, and BER-SE trade-off results discussed in Chapter 5.

\begin{small}
\begin{longtable}{|>{\RaggedRight\arraybackslash}p{0.42\linewidth}|>{\RaggedRight\arraybackslash}p{0.48\linewidth}|}
\caption{MATLAB source files used in the thesis simulation framework}
\label{tab:appendix_matlab_files}\\
\hline
\textbf{File} & \textbf{Role in the simulation} \\
\hline
\endfirsthead

\hline
\textbf{File} & \textbf{Role in the simulation} \\
\hline
\endhead

\path{OTFS_sample_code.m} & Main experiment runner. It loads the configuration, generates OHD and B-OHD pattern tables, runs OFDM-LMMSE, OTFS, OTFS-IM, and random-pattern diagnostics, saves results, and generates figures. \\
\hline
\path{config/config_otfs_im.m} & Defines the numerical simulation setup, including grid size, frame count, SNR range, modulation order, IM block size, random seeds, and reporting options. \\
\hline
\path{OTFS_modulation.m} & Implements baseline OTFS modulation from the delay-Doppler grid to the time-domain transmit signal. \\
\hline
\path{OTFS_demodulation.m} & Implements the receiver-side OTFS demodulation stage and recovers the delay-Doppler-domain observation grid. \\
\hline
\path{OTFS_channel_gen.m} & Generates the discrete multipath delay-Doppler channel taps used in the Monte Carlo simulations. \\
\hline
\path{OTFS_channel_output.m} & Applies the time-varying multipath channel and complex AWGN to the transmitted time-domain signal. \\
\hline
\path{OTFS_mp_detector.m} & Implements the message passing detector used for delay-Doppler symbol detection. \\
\hline
\path{simulation/simulate_baseline_ofdm.m} & Simulates the OFDM-LMMSE reference system used as the conventional multicarrier benchmark. \\
\hline
\path{simulation/simulate_baseline_otfs.m} & Simulates the baseline OTFS system without index modulation. \\
\hline
\path{simulation/simulate_otfs_im.m} & Simulates OTFS-IM using a selected activation-pattern table and records total BER, index BER, symbol BER, and pattern error rate. \\
\hline
\path{simulation/compare_random_patterns.m} & Evaluates random activation-pattern tables as a diagnostic baseline for pattern-selection validation. \\
\hline
\path{pattern_selection/select_patterns_ohd.m} & Selects activation patterns according to the optimized Hamming-distance criterion. \\
\hline
\path{pattern_selection/select_patterns_balanced_ohd.m} & Refines OHD selection by adding a carrier-usage balancing criterion when multiple candidate sets have similar distance properties. \\
\hline
\path{pattern_selection/evaluate_pattern_table.m} & Computes pattern-table metrics, including Hamming-distance statistics and carrier-usage distribution. \\
\hline
\path{reporting/save_simulation_results.m} & Saves the configuration metadata, numerical BER results, pattern information, and generated output files. \\
\hline
\path{reporting/plot_pattern_method_results.m} & Generates per-configuration figures such as BER overview, BER zoom, error decomposition, pattern geometry, and pattern-validation plots. \\
\hline
\path{reporting/plot_se_ber_tradeoff_from_results.m} & Combines saved results from multiple configurations to produce the cross-configuration BER-SE trade-off figures. \\
\hline
\end{longtable}
\end{small}

\subsection*{Appendix B. Reproduction Workflow}

The main numerical results are generated by running \path{OTFS_sample_code.m} once for each selected OTFS-IM configuration. Before each run, the IM block size \(n\) and number of active positions \(k\) are set in \path{config/config_otfs_im.m}. The mathematical definitions of \(b_1\), \(b_2\), spectral efficiency, and the selected \((n,k)\) configurations are already given in Section 5.1, so they are not repeated here.

The basic execution sequence is shown below.

\begin{lstlisting}[language=Matlab,basicstyle=\ttfamily\small,frame=single]
% Step 1: set cfg.n and cfg.k in config/config_otfs_im.m

% Step 2: run the main experiment for this configuration
OTFS_sample_code

% Step 3: after all selected configurations have been saved,
% generate the cross-configuration BER-SE figures
run_plot_se_ber_tradeoff
\end{lstlisting}

For reproducibility, the implementation fixes separate random seeds for the main simulation, the baseline comparison, the OTFS-IM comparison, and the random-pattern diagnostic. This makes it possible to rerun the same configuration while keeping the channel and random-pattern comparison behavior consistent. The main reproducibility-related options are summarized in Table~\ref{tab:appendix_reproducibility_options}.

\begin{table}[H]
\centering
\caption{Main reproducibility and reporting options}
\label{tab:appendix_reproducibility_options}
\footnotesize
\setlength{\tabcolsep}{3pt}
\begin{tabular}{|p{0.46\linewidth}|p{0.10\linewidth}|p{0.34\linewidth}|}
\hline
\textbf{Configuration variable} & \textbf{Value} & \textbf{Purpose} \\
\hline
\texttt{cfg.rng\_seed} & 1 & Main random seed for the experiment runner. \\
\hline
\texttt{cfg.rng\_seed\_baseline} & 11 & Shared seed for OFDM-LMMSE and baseline OTFS comparison. \\
\hline
\texttt{cfg.rng\_seed\_im\_compare} & 22 & Shared seed for OHD and B-OHD OTFS-IM comparison. \\
\hline
\texttt{cfg.rng\_seed\_random\_patterns} & 101 & Seed used for random-pattern table generation. \\
\hline
\texttt{cfg.enable\_random\_compare} & true & Enables the random-pattern diagnostic baseline. \\
\hline
\texttt{cfg.num\_random\_tables} & 5 & Number of random pattern tables used for averaging. \\
\hline
\texttt{cfg.save\_results} & true & Saves numerical results and configuration metadata. \\
\hline
\texttt{cfg.results\_dir} & \texttt{results} & Output directory for saved simulation results. \\
\hline
\end{tabular}
\end{table}

\subsection*{Appendix C. Generated Result Files}

Each run produces configuration-specific result files under the \path{results} directory. The exact numerical results are then used by the reporting functions to generate the figures included in Chapter 5. The most important generated figure categories are summarized in Table~\ref{tab:appendix_generated_outputs}.

\begin{table}[H]
\centering
\caption{Generated result categories}
\label{tab:appendix_generated_outputs}
\begin{tabular}{|p{0.32\linewidth}|p{0.58\linewidth}|}
\hline
\textbf{Result category} & \textbf{Purpose} \\
\hline
BER overview & Compares OFDM-LMMSE, baseline OTFS, OHD OTFS-IM, and B-OHD OTFS-IM for one \((n,k)\) configuration. \\
\hline
BER zoom & Focuses on the OTFS and OTFS-IM curves in the medium-to-high-SNR region. \\
\hline
Error decomposition & Separates OTFS-IM errors into index BER, symbol BER, and pattern error rate. \\
\hline
Pattern geometry & Shows selected activation patterns, Hamming-distance structure, and carrier-usage distribution. \\
\hline
Random-pattern validation & Compares OHD and B-OHD pattern selection with random pattern tables. \\
\hline
Configuration BER summary & Summarizes BER across several \((n,k)\) configurations at selected SNR points. \\
\hline
BER-SE trade-off & Relates BER performance to spectral efficiency across all simulated configurations. \\
\hline
\end{tabular}
\end{table}

No additional appendix figures are required because the representative figures are already included in Chapter 5 at the point where they are analyzed. Adding the same result figures again in the appendix would duplicate the main discussion. The appendix is therefore kept as an implementation and reproducibility reference.
